\section{Communication summary (reconstructed project-cycle transcript)}
\label{sec:app_communication_summary}

This appendix records a reconstructed full communication summary for the current thrust-analysis cycle.
The source material is the project artifact trail (plans, note rewrites, workflow outputs, and validation products), not a verbatim message log.
The objective is to preserve the decision chain for internal reproducibility and method auditing.

\subsection{Scope and reconstruction sources}

Scope covered: the thrust Iterative Bayesian Unfolding cycle from code review through preliminary-result packaging and note rewrites (Sections 1--8), including the AI-agent collaboration framing.

Primary reconstruction sources:
\begin{itemize}
\item dated planning artifacts under \texttt{Plans/} (2026-02-13 to 2026-02-18);
\item executable workflow outputs in \texttt{Studies/PreliminaryResult/}, with run directory \texttt{2026-03-05\_thrust\_iter5/};
\item note section revision trail in \texttt{AnalysisNote/sections/*.tex};
\item workflow scripts in \texttt{Unfolding/scripts/} and \texttt{Studies/*/scripts/}.
\end{itemize}

\subsection{Reconstructed communication timeline}

\begin{enumerate}
\item \textbf{Initial objective setting (2026-02-13).} The collaboration established a thrust-first execution strategy with Iterative Bayesian Unfolding, complete code review, and a JHEP-style internal-note structure.
\item \textbf{Code-audit phase.} Critical and high-impact defects were identified and fixed (tree defaults, closure mode requirements, return-code handling, systematics path robustness), with findings tracked in \texttt{Plans/2026-02-13\_code\_review\_findings.md}.
\item \textbf{1994 workflow focus.} Execution direction was corrected to run the full 1994 derivation and unfolding chain rather than legacy partial-year paths.
\item \textbf{Nominal-regularization policy convergence.} The cycle converged on nominal $N_{\mathrm{iter}}=5$ with 4/6 as regularization systematic envelope for this thrust preliminary pass.
\item \textbf{Systematics completion requirement.} The scope was expanded to include full experimental and theory branches in the same run-owned output package, with hadronic correction retained in nominal processing.
\item \textbf{Result-style iteration loop.} Multiple rounds of visual QA enforced consistent pad geometry, ratio-panel readability, uncertainty normalization conventions, and legend/text placement for publication-grade interpretation.
\item \textbf{ALEPH comparison integration.} Final-result plotting was upgraded to include ratio to ALEPH 2004 in matched binning with uncertainty propagation via ROOT histogram operations.
\item \textbf{Uncertainty decomposition policy.} The workflow migrated from envelope-only views to explicit component-level plots, including signed shifts, grouped uncertainty structure, and central-window interpretation.
\item \textbf{Section-by-section rewrite cycle (Sections 1--5 first).} The note was expanded with stronger physics motivation, detector/data context, and reconstruction validation text connected directly to thrust sensitivity.
\item \textbf{Sections 5+ hardening.} Methods, systematics, and results sections were restructured for executable traceability (script entrypoints, output paths, and release artifacts).
\item \textbf{AI-agent narrative formalization.} A dedicated section captured operating model, controls, and measurable scorecard, followed by a final chapter logging persistent collaboration and analysis challenges.
\item \textbf{Physics-tone refinement.}\\
Introductory framing was strengthened around precision-QCD motivation,\\
LEP data quality, and validation goals.
\item \textbf{External-citation expansion.}\\
AI-for-science and agent-workflow references were expanded, then extended with requested recent case studies.
\item \textbf{Reproducibility emphasis.} The communication repeatedly enforced path-level provenance for figures and summary products, with explicit checks and output validation scripts.
\item \textbf{Current-cycle layout request.} Results narrative was requested to present uncertainty decomposition before the final unfolded comparison, and to add a final-result total correlation matrix view.
\item \textbf{Current-cycle documentation request.} A complete communication appendix was requested for community context, with highlighted key decisions and rationale.
\end{enumerate}

\subsection{Key exchanges and decisions}

\newcommand{\decisionbox}[3]{%
\noindent\fbox{%
\parbox{0.97\textwidth}{%
\textbf{Decision:} #1\\
\textbf{Rationale:} #2\\
\textbf{Impact:} #3
}}%
}

\decisionbox
{Use Iterative Bayesian Unfolding as the baseline correction method for this cycle.}
{Fast iteration without model retraining was required for rapid experimental cross-check loops.}
{Workflow scripts, method text, and uncertainty propagation were all aligned to IBU.}

\decisionbox
{Set nominal regularization to $N_{\mathrm{iter}}=5$ with 4/6 as regularization systematic.}
{Iteration-scan diagnostics and practical stability/throughput tradeoffs favored this operating point.}
{All nominal and systematic branches, note text, and final plots were normalized to the same policy.}

\decisionbox
{Run the full 1994 chain as the maintained baseline dataset year.}
{The analysis state required consistency in data/MC derivation and response inputs for one complete production year.}
{Output provenance and note references were consolidated under run-owned 1994 directories.}

\decisionbox
{Adopt a strict figure-style contract for result plots (connected pads, controlled ratio range, readable labels).}
{Visual interpretation quality was treated as analysis-relevant, not cosmetic.}
{Several plotting scripts were iteratively corrected until style and uncertainty semantics stabilized.}

\decisionbox
{Promote component-level uncertainty views (statistical, unfolding, experimental, theory) over envelope-only summaries.}
{Readers need source-resolved behavior rather than only combined envelopes.}
{Results section now reports component behavior explicitly and ties interpretation to central-window metrics.}

\decisionbox
{Keep ``no-neutrals'' as a cross-check branch and separate it from nominal total uncertainty.}
{It is a stress test rather than a baseline uncertainty source.}
{Cross-check visibility is preserved while nominal total remains policy-consistent.}

\decisionbox
{Introduce explicit AI-agent operating-model and challenge-tracking sections.}
{The internal-note objective includes evaluating whether an agent can execute research workflows under supervision.}
{The note now contains controls, measurable scorecard, and persistent challenge logging.}

\decisionbox
{Expand citations to include recent AI+scientist and theory-assistant case studies.}
{Program motivation required concrete external evidence beyond generic AI discussion.}
{Introduction and AI-agent sections now cite both broad and recent workflow-relevant studies.}

\decisionbox
{Reorder results narrative to show uncertainty structure before final unfolded comparison.}
{Interpretation of the final spectrum is clearer once uncertainty sources are established first.}
{Section flow now progresses from component structure to correlated final products to central comparison.}

\decisionbox
{Add a final-result total correlation matrix plot from \texttt{corr\_total.csv}.}
{Community-facing usage requires direct visibility of correlated uncertainties after full nominal combination.}
{Release package now includes a dedicated final total-correlation figure and script-level provenance.}

\subsection{Communication-pattern summary}

The reconstructed cycle indicates a high-frequency specification--execution loop with strong emphasis on:
\begin{itemize}
\item explicit physics-policy locking before broader text expansion;
\item figure-level interpretability as a first-class requirement;
\item command/path-level provenance for every released artifact;
\item incremental correction and revalidation rather than large unreviewed batches.
\end{itemize}

The dominant friction points were precision formatting constraints, rapid scope evolution, and synchronization between code outputs and note narrative. These patterns are consistent with the persistent-challenge tracker in Sec.~\ref{sec:challenges}.
